\section{Certificates, relaxation comparisons, and latent separators}
\label{sec:certification}

A useful design bound does not require an optimization run to terminate
successfully: a feasible schedule and a separately checked upper witness
suffice. This section gives such witnesses, compares them with established
virtual-noise relaxations, and constructs an alternative exact representation
for cases where a long calendar window is expensive. Throughout this section
$\calF$ is nonempty and finite, and $J_0\succ0$ unless stated otherwise; this
stronger prior assumption ensures that every logarithm and inverse used in
the certificates is defined.

\subsection{Prior-aware support certificates}
\label{subsec:local-certificates}
Suppose a graph represents exactly $\calF$, a path $P_S$ represents schedule
$S$, and the local increments $Q_e\succeq0$ satisfy
\[
 J_L(S)=J_0+\sum_{e\in P_S}Q_e,\qquad
 J(S)\preceq J_0+\frac{1}{1-\delta}\sum_{e\in P_S}Q_e,
 \qquad 0\le\delta<1.
\]
The second inequality is the data-information consequence of
Lemma~\ref{lem:residual-transfer}; it retains the prior without inflation.
For scalar, complete-block, general-decay, and partial-packet models,
Section~\ref{sec:locality} supplies the corresponding sufficient $\delta$.

\begin{theorem}[Support certificate]\label{thm:local-certificate}
For any $N\succ0$, let
\begin{align}
 c_N&=\log\det N-p+\tr(N^{-1}J_0),\nonumber\\
 a_e&=\frac{\tr(N^{-1}Q_e)}{1-\delta},\qquad
 U_N=c_N+\max_{S\in\calF}\sum_{e\in P_S}a_e.
 \label{eq:local-certificate}
\end{align}
Then $\max_{S\in\calF}\log\det J(S)\le U_N$.
If $S_0\in\calF$ and $\ell\le\log\det J(S_0)$, an upper enclosure
$U\ge U_N$ certifies D-efficiency at least $\exp\{(\ell-U)/p\}$
relative to the true optimum.
\end{theorem}
\begin{proof}
For any $X\succ0$, concavity gives
$\log\det X\le\log\det N-p+\tr(N^{-1}X)$.
Apply monotonicity of log determinant to the prior-aware upper matrix,
then this tangent inequality, and maximize the additive term. The efficiency
statement follows by exponentiating the difference between the incumbent
lower bound and the optimum upper bound.
\end{proof}

For a weighted-trace objective with $W\succeq0$, the same matrix bound gives
\[
 \max_{S\in\calF}\tr(WJ(S))\le \tr(WJ_0)
 +\max_{S\in\calF}\sum_{e\in P_S}\frac{\tr(WQ_e)}{1-\delta}.
\]
This requires no logarithm or tangent reference; the partial-packet trace
certificates use it. Dividing an exact feasible lower
value by a positive upper value certifies relative performance; a zero upper
value proves that every feasible value is zero.

Neither $N$ nor $J_0+(N-J_0)(1-\delta)$ needs to be the information of a
feasible schedule. For example, optimizing log determinant over a convex
mixture of path matrices can propose $N$, and an exchange heuristic can
propose $S_0$; the mixture supplies no feasible-experiment lower bound. The
support maximum in \eqref{eq:local-certificate} must range over all feasible
paths, not only over paths discovered by an optimizer. Exact cardinality,
mandatory/forbidden choices, and minimum gap are enforced by the count, mask,
and cooldown graph of Section~\ref{subsec:calendar-graph}; a spacing
constraint remains in the graph even if its required cooldown exceeds the
information window.

The theorem permits real witnesses. For the exact rational checker, choose
a rational symmetric reference $N$ (or its rational inverse), for example by
rounding a numerical proposal, and verify $N\succ0$ by exact positive pivots
or leading principal minors; a failed check rejects the witness. With rational
model data, the information of $S_0$, every $Q_e$, and $N^{-1}$ can then be
recomputed by rational elimination. The weighted-trace checker likewise
requires its prescribed weight $W$ to be rational and verifies $W\succeq0$.
Because rational model data need not make the analytic locality bound
rational, every score uses a certified rational enclosure
$\delta\le\widehat\delta<1$ in place of $\delta$. This enlargement preserves
the upper matrix bound because the increments are PSD. Rational envelopes or
outward interval evaluation of the locality formulas give such enclosures;
square-root endpoints can be verified by squaring. If the enclosure is not
below one, refine it, increase the memory window and recompute the bound, or
reject that certificate.
No numerical reference, solver status, or reported gradient is used as a
proof. Appendix~\ref{app:certificate-arithmetic} gives complete outward score
and logarithm formulas. If $a_e$ is replaced by
$\lceil G_s a_e\rceil/G_s$ and every path has at most $k$ nonzero selection
arcs, the support increase is at most $k/G_s$ (and is zero when $k=0$).
More general integer interval scores remain valid, with their actual excess
explicitly charged. The rational certificate proves a statement about the
encoded inputs. Exact arithmetic does not remove uncertainty in the physical
covariance or the sensitivity construction; that requires additional
assumptions.

Logdet support conditions for finite information hulls are classical
\citep{HarmanTrnovska2009}, and correlated-design simplicial decomposition is
already developed by \citet[Section 4.4]{Hainy2025}. The role of
Theorem~\ref{thm:local-certificate} is to attach an explicit uniform
selected-covariance error bound to this established support machinery.

\subsection{The virtual-noise comparator and exact continuous bounds}
\label{subsec:virtual-noise}
Here we index scalar candidate responses by $i=1,\ldots,n$; packet
selection ties their coordinates. Let
\[
 D=\diag(a_i)\succ0,\qquad B=R-D\succeq0,
 \qquad Z(z)=\diag(z_i/a_i),\quad z\in[0,1]^n.
\]
The established covariance-split extension is
\begin{equation}\label{eq:vn-resolvent}
 V(z)=(I+BZ(z))^{-1}F,\qquad
 J_D(z)=J_0+F^T Z(z)V(z),\qquad \varphi_D(z)=\log\det J_D(z).
\end{equation}
These formulas do not invert $Z(z)$ and therefore apply at zero weights.
The matrix $I+BZ$ is invertible: Sylvester's determinant identity reduces
its determinant to that of $I+B^{1/2}ZB^{1/2}\succ0$.

\begin{proposition}[Established virtual-noise identities]
\label{prop:vn}
For $T=\{i:z_i>0\}$,
\begin{equation}\label{eq:vn-covariance}
 J_D(z)=J_0+F_T^T
 \left[R_{TT}+\diag\left(a_i\frac{1-z_i}{z_i}:i\in T\right)\right]^{-1}F_T.
\end{equation}
The empty-support value is $J_0$. The map $J_D$ is matrix concave on the
cube, $\varphi_D$ is concave and differentiable there by extension to a
neighborhood, and
\begin{equation}\label{eq:vn-gradient}
 \frac{\partial J_D}{\partial z_i}=a_i^{-1}V_i^TV_i,\qquad
 g_i(z):=\frac{\partial\varphi_D}{\partial z_i}
       =a_i^{-1}V_iJ_D(z)^{-1}V_i^T.
\end{equation}
For binary $z=\ind_S$, $J_D(z)=J(S)$, independently of $D$.
\end{proposition}
\begin{proof}
The nonzero rows of $ZV$ solve
$(B_{TT}+Z_{TT}^{-1})(ZV)_T=F_T$, proving
\eqref{eq:vn-covariance}. For $B\succ0$,
\[
 Z(I+BZ)^{-1}=B^{-1}-B^{-1}(B^{-1}+Z)^{-1}B^{-1}
\]
is matrix concave because matrix inversion is convex. Approximate
$B\succeq0$ by $B+\epsilon I$ and pass to the continuous resolvent limit;
this proves concavity also on the split boundary. Monotone concave log
determinant preserves concavity. Differentiating the resolvent gives
$d\{Z(I+BZ)^{-1}\}=(I+ZB)^{-1}(dZ)(I+BZ)^{-1}$,
which proves \eqref{eq:vn-gradient}, including the cube boundary.
Binary substitution into \eqref{eq:vn-covariance} proves exactness.
\end{proof}

The scalar choice $D=aI$ is the Liu relaxation
\citep[Section II.B]{Liu2016}. Its equality with the convex virtual-noise
design of \citet{Pazman2022}, under the measure change $\xi_i=z_i/k$ for
$k>0$, was already proved by \citet[Proposition 3]{Hainy2025}, including
zero weights. Their modified variance uses
$D=\widetilde\kappa\diag(R_{ii})$ with
$\widetilde\kappa\le\lambda_{\min}(\diag(R_{ii})^{-1/2}
R\diag(R_{ii})^{-1/2})$. This is a subfamily of the diagonal splits above;
for constant $R_{ii}=v$, it is the scalar family $a=v\widetilde\kappa$.
The equivalence and the concavity are prior results, not contributions of
this paper.

Let $\mathcal Z$ be a specified compact convex relaxation satisfying
\[
 \operatorname{conv}\{\ind_S:S\in\calF\}
 \subseteq\mathcal Z\subseteq[0,1]^n.
\]
A rational query $z$ in the cube gives
\begin{equation}\label{eq:vn-continuous-certificate}
 \max_{x\in\mathcal Z}\varphi_D(x)
 \le \overline{\log\det J_D(z)}-g(z)^Tz
       +\max_{x\in\mathcal Z}g(z)^Tx.
\end{equation}
Here the overline is a rigorous logarithm upper enclosure, and the last
quantity must be computed or bounded globally. For
$\mathcal Z_k=\{x\in[0,1]^n:\sum_i x_i=k\}$ it is the sum of the $k$
largest rational gradient coordinates. If $\mathcal Z$ is the true indicator
hull with spacing constraints, linear optimization uses its feasible graph
instead. At a verified feasible $z\in\mathcal Z$,
$\underline{\log\det J_D(z)}$ is a lower bound on the continuous optimum.
An infeasible query still supplies a tangent upper bound but no such lower
bound. Thus continuous optimization error and the discrete relaxation gap
can be measured separately. A strict comparison between a continuous lower
bound and a competing integer-objective upper bound does not require knowing
the integer optimum.

Exact dense elimination suffices to check all these expressions. For scalar
AR(1)-plus-nugget covariance, Appendix~\ref{app:dense-oracle} derives
covariance-filter and reverse-adjoint formulas that are linear in $n$, and
exact tridiagonal split checks; these implement the same relaxation.
Floating-point precision-space solves near unit latent correlation, and a
first implementation based on subtracting RTS smoothed means, can produce
invalid tangents; the appendix explains the two failure mechanisms.
Only recomputed rational witnesses make \eqref{eq:vn-continuous-certificate}
a rigorous bound.

\subsection{A certificate against every diagonal split}
\label{subsec:all-split}
A virtual-noise relaxation can appear weak merely because it was optimized
incompletely. The following comparison rules out both incomplete continuous
solution and tuning within the diagonal-split family. It does not compare
complete integer branch-and-bound methods.

Fix $z^0\in\mathcal Z$ and set $T=\{i:z_i^0>0\}$,
$h_i=(1-z_i^0)/z_i^0$ for $i\in T$. For every positive vector $a$, define
\[
 S(a)=R_{TT}+\diag(h_i a_i:i\in T),\qquad
 \phi(a)=\log\det\{J_0+F_T^T S(a)^{-1}F_T\}.
\]
This fixed-point value exists even if $R-\diag(a)$ is not PSD.
For $T=\varnothing$, use $J=J_0$, $g=0$, and omit all empty matrix terms.

\begin{proposition}[All-split lower witness]\label{prop:diagonal-barrier}
Choose any $a^0>0$, put $V=S(a^0)^{-1}F_T$, $J=J_0+F_T^TV$,
and let
\[
 g_i=-h_i V_iJ^{-1}V_i^T\quad(i\in T),\qquad g_i=0\quad(i\notin T).
\]
For every $Y\succeq0$ with $\diag(Y)\ge-g$, the number
\begin{equation}\label{eq:diagonal-barrier}
 L_{\rm split}=\phi(a^0)-g^Ta^0-\tr(RY)
\end{equation}
satisfies $\varphi_D(z^0)\ge L_{\rm split}$ for every admissible
$D=\diag(a)\succ0$, $R-D\succeq0$. Consequently it lower-bounds every
such continuous optimum and
\begin{equation}\label{eq:pointwise-split-barrier}
 \sup_{z\in\mathcal Z}\inf_{D}\varphi_D(z)\ge L_{\rm split}.
\end{equation}
Every relaxation consisting of $z\in\mathcal Z$ and any collection of
globally valid affine upper supports $t\le l_D(z)$ from these splits also
has optimum at least $L_{\rm split}$.
\end{proposition}
\begin{proof}
Put $C=F_TJ_0^{-1}F_T^T\succeq0$. The determinant lemma gives
$\phi(a)=\log\det J_0+\log\det(S(a)+C)-\log\det S(a)$.
This is a classical convex function of $S$ \citep{KimKim2006,Kim2019}.
For completeness, with $A=S^{-1}\succeq B_1=(S+C)^{-1}$, its Hessian
in a symmetric direction $H$ is
$\tr(AHAH)-\tr(B_1HB_1H)\ge0$:
$A\otimes A-B_1\otimes B_1=(A-B_1)\otimes A+B_1\otimes(A-B_1)\succeq0$.
Affine substitution proves convexity in $a$, with gradient $g$ as displayed.
For $w=-g\ge0$, admissibility and the dual conditions imply
\[
 w^Ta\le\tr(DY)\le\tr(RY).
\]
The convex lower tangent $\phi(a)\ge\phi(a^0)+g^T(a-a^0)$ therefore
proves \eqref{eq:diagonal-barrier}. The same point $z^0$ works for all $D$,
which proves \eqref{eq:pointwise-split-barrier} without a minimax interchange.
For any globally valid affine upper support,
$l_D(z^0)\ge\varphi_D(z^0)\ge L_{\rm split}$. Hence $(z^0,L_{\rm split})$
remains feasible in their intersection.
\end{proof}

The proof uses established logdet convexity and weak semidefinite duality;
its useful output is a finite checkable witness against an entire family.
For rational $a^0$, all entries except the logarithm are rational. Any
rational factor $B_0$ gives a valid dual matrix by construction:
\begin{equation}\label{eq:dual-repair}
 Y=B_0B_0^T+\diag\left(\max\{0,-g_i-\sum_j(B_0)_{ij}^2\}\right).
\end{equation}
A floating-point SDP solution can propose $B_0$ but is not assumed feasible.
The lower logarithm enclosure is used in \eqref{eq:diagonal-barrier}.
The reference $a^0$ need not be admissible, and neither strong duality
nor optimality of this reference or dual factor is required.

There is a simpler scalar witness. For fixed $z^0$, increasing $a$ increases
$R_{TT}+a\diag(h)$, so $\varphi_{aI}(z^0)$ is nonincreasing. Therefore
\begin{equation}\label{eq:scalar-barrier}
 \bar a\ge\lambda_{\min}(R)\quad\Longrightarrow\quad
 \max_{z\in\mathcal Z}\varphi_{aI}(z)
 \ge\varphi_{\bar aI}(z^0)\quad
 (0<a\le\lambda_{\min}(R)).
\end{equation}
The right side is evaluated using \eqref{eq:vn-covariance}, not by an
inadmissible concavity assertion. Any nonzero rational $v$ supplies
$\bar a=v^TRv/(v^Tv)\ge\lambda_{\min}(R)$; alternatively an exact nonpositive
leading pivot of $R-\bar aI$ proves the same inequality. Appendix~\ref{app:dense-oracle}
gives the smaller tridiagonal witness for a stationary latent chain.

\begin{example}[A persistent gap]\label{ex:split-barrier}
Let $R=\left(\begin{smallmatrix}2&1\\1&2\end{smallmatrix}\right)$,
$F=(1,-1)^T$, $J_0=1$, and $k=1$. Every feasible schedule has information
$3/2$. At $z^0=(1/2,1/2)$ and $a^0=(1,1)$, $J=2$,
$g=(-1/8,-1/8)$, and
$Y=\frac18\left(\begin{smallmatrix}1&-1\\-1&1\end{smallmatrix}\right)$
has $\diag Y=-g$ and $\tr(RY)=-g^Ta^0=1/4$.
Thus $L_{\rm split}=\log2$, leaving a gap at least $\log(4/3)$ over the
integer optimum for every diagonal split, including strict splits and
pointwise split optimization. All affine upper cuts in the stated family
retain this witness. Branching, additional restrictions excluding $z^0$,
or cuts outside this family can remove it. This elementary example is not
a novelty claim for the relaxation or the convexity argument.
\end{example}

\subsection{One shared schedule over finite scenarios}
\label{subsec:robust-certificate}
Let $s=1,\ldots,q$ index a declared finite set of known covariance and
sensitivity models, all with information dimension $p$, positive definite
priors, and the same nonempty feasible schedule family $\calF$. Define
\begin{align}
 f_s(S)&=\log\det J_s(S),\qquad
 b_s^*=\max_{S\in\calF}f_s(S),\nonumber\\
 g_c(S)&=\min_s(f_s(S)-c_s),\qquad
 g_*(S)=\min_s(f_s(S)-b_s^*),\qquad
 E_*(S)=\exp(g_*(S)/p).\label{eq:robust-objectives}
\end{align}
Raw maximin sets $c_s=0$. Standardized D-efficiency compares each scenario
with its own best attainable information. Under a scenario-dependent
invertible congruence $J_s\mapsto T_sJ_sT_s^T$, both $f_s$ and $b_s^*$
gain $2\log|\det T_s|$, so $g_*$ is invariant if priors transform too.
The maximizing schedules of raw maximin are invariant under a common
transformation, but need not be under unrelated transformations.

Efficiency normalization and its support-based optimization are established:
\citet[Section 3.2 of the preprint]{Burclova2016} computes individual
criterion optima before maximizing minimum efficiency by affine supports,
and \citet{Maus2010} studies maximin relative efficiency under uncertain
AR(1) correlation. Robust kinetic design and scenario exchange also have
direct precedents \citep{Duarte2018,WangYue2026}; the latter preprint uses
independent-trial sampling information. \citet{Chowdhary2025} develops
budget-constrained robust correlated-noise sensor design through stochastic
policy optimization. Our contribution is an explicit certificate for
one common discrete schedule that rigorously propagates the uncertainty in
its optimal standardizers.

Suppose the common path graph supplies $Q_{s,e}\succeq0$ and $\delta_s<1$
with $J_s(S)\preceq J_{0,s}+\sum_{e\in P_S}Q_{s,e}/(1-\delta_s)$.
Choose any $N_s\succ0$ and exactly nonnegative $\lambda_s$ with
$\sum_s\lambda_s=1$. Then
\begin{equation}\label{eq:robust-support}
 \max_{S\in\calF}g_c(S)\le
 \sum_s\lambda_s[\log\det N_s-p+\tr(N_s^{-1}J_{0,s})-c_s]
 +\max_{S\in\calF}\sum_{e\in P_S}\sum_s
 \frac{\lambda_s\tr(N_s^{-1}Q_{s,e})}{1-\delta_s}.
\end{equation}
Indeed, the minimum is at most a convex average, each scenario obeys
Theorem~\ref{thm:local-certificate}, and one common path prices the combined
scores. Separate scenario support maxima also give a valid bound but may
be strictly weaker. The rational checker uses rational SPD references $N_s$
and certified rational enclosures $\delta_s\le\widehat\delta_s<1$ in all
scenario scores, as above. It also uses rational simplex weights: round a
numerical proposal to rationals, clip negative entries, and normalize exactly
if positive mass remains; otherwise use a specified rational simplex point
or reject. Approximately normalized or negative weights do not satisfy the
proof.

\begin{proposition}[Unknown standardizers]\label{prop:robust-normalizers}
Suppose verified bounds $\ell_s\le b_s^*\le U_s$ are available.
For a common feasible $S_0$, let
\begin{equation}\label{eq:robust-certified-interval}
 L=\min_s[\underline{f_s(S_0)}-U_s].
\end{equation}
In \eqref{eq:robust-support}, replace $c_s$ by $\ell_s$ and round every
term outward to obtain an upper bound $U'$. Then
\[
 L\le g_*(S_0)\le\max_{S\in\calF}g_*(S)\le U:=\min(U',0),
 \quad E_*(S_0)\ge e^{L/p},\quad
 \frac{E_*(S_0)}{\max_{S\in\calF}E_*(S)}\ge e^{(L-U)/p}.
\]
\end{proposition}
\begin{proof}
Pointwise,
$\min_s(f_s(S)-U_s)\le g_*(S)\le\min_s(f_s(S)-\ell_s)$.
The left inequality gives the incumbent lower bound; the right and
\eqref{eq:robust-support} give the upper bound. Since $f_s(S)\le b_s^*$,
$g_*(S)\le0$. Exponentiation gives both efficiency conclusions.
\end{proof}
The individual lower bounds $\ell_s$ can come from different schedules, but
each must be feasible for the same family defining $b_s^*$. A reference
optimum over a larger family gives an upper bound, not automatically a valid
lower bound. The rational log interval is the certificate; floating-point
exponentials are only display summaries. These statements certify exactly
the listed scenarios and do not cover an intervening continuous parameter
region.

\begin{corollary}[Fixed-scenario approximation]\label{cor:robust-fptas}
Under the fixed rational memory promises and feasibility conditions of
Theorem~\ref{thm:trace-fptas}, with $q$ and $p$ fixed, standardized
D-efficiency has an FPTAS. This is a direct consequence of the spectral
approximation set, not a runtime guarantee for the practical mixture solver.
\end{corollary}
\begin{proof}
Take one graph with the largest required memory. Apply
Theorem~\ref{thm:dag-spectral-set} to the block-diagonal additive matrix
$\diag(J_{L,1}(S),\ldots,J_{L,q}(S))$ of fixed order $qp$.
If every memory error is at most $\delta$, every target schedule has a
single returned representative satisfying, in every scenario,
\begin{equation}\label{eq:robust-cover-factor}
 J_s(\widehat S)\succeq aJ_s(S),\qquad
 a=\frac{(1-\eta)(1-\delta)}{1+\delta}.
\end{equation}
This follows by transferring true information to local information, using
the common spectral cover, and transferring back. For the returned set
$\mathcal C$, compute the positive rational numbers
$d_s=\max_{S\in\mathcal C}\det J_s(S)$ and choose $\widehat S\in\mathcal C$
maximizing $\min_s\det J_s(S)/d_s$. If
$D_s^*=\max_{S\in\calF}\det J_s(S)$, then
$a^pD_s^*\le d_s\le D_s^*$. Consequently the efficiency $E_{\mathcal C}$
normalized by $d_s$ satisfies
$E_*(S)\le E_{\mathcal C}(S)\le E_*(S)/a$.
A representative $S'$ of a robust optimum has
$E_*(S')\ge a\max_S E_*(S)$, whence
\[
 E_*(\widehat S)\ge a E_{\mathcal C}(\widehat S)
 \ge a E_{\mathcal C}(S')\ge a^2\max_S E_*(S).
\]
Choose $\eta=\epsilon/4$ and $\delta\le\epsilon/8$, or full history
with $\delta=0$. Then $a\ge1-\eta-2\delta\ge1-\epsilon/2$ and
$a^2\ge1-\epsilon$. The graph and output set have polynomial size under
the stated fixed promises; exact true determinants and the displayed
rational comparisons have polynomial bit complexity. No logarithmic or
root comparison oracle is needed. The polynomial degree depends on $qp$.
\end{proof}
With known positive rational determinant normalizers, maximizing the
corresponding rational determinant ratios over $\mathcal C$ loses only the
single factor $a$ in determinant-root efficiency. Unknown normalizers cost a
second factor in the same-set procedure above. Appendix~\ref{app:robust-witnesses}
gives a sharp tie example and a common-mixture integrality gap.

\subsection{Exact latent-separator representations}
\label{subsec:separators}
When a fixed physical covariance is sampled on finer grids, the per-step
contraction can approach one and the sufficient calendar window can become
large. An alternative keeps selected latent states as nuisance variables
and conditions short blocks on them exactly, without a truncation parameter.
For this construction let
\begin{equation}\label{eq:separator-model}
 Y=F\theta+X+v,\quad X\sim N(0,K),\quad
 v\sim N(0,rI),\quad K\succ0,\ r>0,
\end{equation}
where $X$ is a scalar Gaussian Markov chain and $X,v$ are independent.
Unlike the singular-state locality theorem, the inverse formulas here require
$K\succ0$; independent or zero-latent cases can be evaluated directly.
For an anchor set $A$, write $U=X_A$, $K_A=K_{AA}$, and
\begin{equation}\label{eq:separator-conditional}
 H_A=K_{:A}K_A^{-1},\qquad
 D_A=K-K_{:A}K_A^{-1}K_{A:}+rI\succ0.
\end{equation}
With $A=\varnothing$, omit the anchor terms and put $D_A=R=K+rI$.
Conditioning on anchors separates the intervening Markov intervals.
Each anchor-time observation is assigned to exactly one neighboring block;
its conditional latent residual is zero and its observation noise remains
independent. Thus $D_A$ is block diagonal over a disjoint partition
$I_1,\ldots,I_m$. For example, partition into consecutive blocks of at most
$b$ times and anchor every block's right endpoint except the final time.
Each block then has at most two adjacent anchors.

For any schedule define the augmented information
\begin{equation}\label{eq:separator-augmented}
 M_A(S)=\diag(J_0,K_A^{-1})+
 [F_S,H_{A,S}]^T(D_A)_{SS}^{-1}[F_S,H_{A,S}].
\end{equation}
It is the fixed prior plus one PSD atom for each within-block pattern
$S\cap I_j$. Both its full matrix and its nuisance principal block are SPD.
Woodbury's identity gives exactly
\begin{equation}\label{eq:separator-schur}
 \mathcal S_\theta(M_A(S))=J_0+F_S^TR_{SS}^{-1}F_S=J(S),
\end{equation}
where $\mathcal S_\theta$ eliminates all nuisance coordinates. The anchors
are latent nuisance variables with prior precision $K_A^{-1}$; they are
not additional observations available at no cost.

For a generic SPD matrix
$M=\left(\begin{smallmatrix}A_0&B_0\\B_0^T&C_0\end{smallmatrix}\right)$
with parameter block of order $p$, put
\[
 J=A_0-B_0C_0^{-1}B_0^T,\quad G_*=-C_0^{-1}B_0^T,
 \quad E_*=[I_p;G_*],\quad f(M)=\log\det J.
\]
Completing the square gives, for every nuisance matrix $G$,
\begin{equation}\label{eq:schur-variational}
 [I_p;G]^TM[I_p;G]
 =J+(G-G_*)^TC_0(G-G_*)\succeq J.
\end{equation}
This identity proves matrix concavity and monotonicity of the Schur map:
its quadratic form in any parameter vector is the minimum over nuisance
vectors of a jointly linear-in-$M$ quadratic form. Thus $f$ is concave.
Differentiation at the minimizing $G_*$ gives
\begin{equation}\label{eq:separator-gradient}
 \nabla f(M)=E_*J^{-1}E_*^T\succeq0,
 \qquad \rank\nabla f(M)=p,\quad
 \tr(\nabla f(M)M)=p.
\end{equation}
The rank remains $p$ although the nuisance dimension grows.

\begin{theorem}[Separator support certificate]\label{thm:separator-certificate}
For any $W\succ0$ of order $p$ and any $|A|\times p$ matrix $G$,
\begin{align}
 \max_{S\in\calF}\log\det J(S)
 &\le -\log\det W-p+\tr(WJ_0)+\tr(WG^TK_A^{-1}G)
 \nonumber\\[-2pt]
 &\quad+\max_{S\in\calF}\sum_{j=1}^m
 \tr\left[W(F_{S_j}+H_{A,S_j}G)^T
 (D_A)_{S_jS_j}^{-1}(F_{S_j}+H_{A,S_j}G)\right],
 \label{eq:separator-certificate}
\end{align}
where $S_j=S\cap I_j$. Empty patterns contribute zero and empty anchor
terms are omitted. If $\calF=\{S:|S|=k\}$, the support term can be priced
by enumerating at most $2^b$ patterns per block and an exact count dynamic
program, using $O(m(k+1)2^b)$ transitions after pattern scores are available.
\end{theorem}
\begin{proof}
Apply \eqref{eq:schur-variational} to \eqref{eq:separator-augmented},
then use $\log\det X\le-\log\det W-p+\tr(WX)$.
Block diagonality of $D_A$ separates the pattern scores and the fixed
anchor/prior term is added once. For block $j$, let $v_j(c)$ be the largest
score among its patterns of size $c$ (or $-\infty$ if none).
Starting with $V_0(0)=0$ and other states $-\infty$, use
$V_j(t)=\max_c\{V_{j-1}(t-c)+v_j(c)\}$. Its final state $V_m(k)$ is the
support maximum. Pattern enumeration plus these transitions satisfies the
stated bound.
\end{proof}
Local forbidden or mandatory observations can be enforced while enumerating
patterns. Cross-block constraints such as minimum spacing require additional
boundary states; expected-count constraints on independent block mixtures
are not a substitute for the hull of complete feasible schedules.
Any rationally rounded $G$ is valid in \eqref{eq:separator-certificate};
only $W\succ0$ needs a definiteness check. Rational conditional covariances,
integer interval pattern scores, and the logarithm bounds of
Appendix~\ref{app:certificate-arithmetic} make this an exact certificate.
At a mixture matrix $M$, choosing $W=J^{-1}$ and $G=G_*$ recovers its ordinary
upper tangent. Arbitrary witnesses need not solve these stationarity equations.

\begin{theorem}[Nested-anchor hierarchy]\label{thm:separator-hierarchy}
For the same model and feasible family, define the mathematical hull optimum
\[
 U_A=\max_{\alpha\in\Delta(\calF)}
 \log\det\mathcal S_\theta\left(\sum_{S\in\calF}\alpha_SM_A(S)\right).
\]
For nested anchor sets $A\subseteq B$,
\begin{equation}\label{eq:separator-order}
 \max_{S\in\calF}\log\det J(S)\le U_A\le U_B.
\end{equation}
Thus removing anchors and merging conditional blocks tightens these hull
bounds. The statement compares complete hull optima, not unfinished
numerical upper witnesses or nonnested partitions.
\end{theorem}
\begin{proof}
Order $M_B$ by $(\theta,X_A,X_{B\setminus A})$ and let
$\mathcal S_{\theta,A}$ eliminate only $X_{B\setminus A}$.
The exact identity
\begin{equation}\label{eq:separator-cross-identity}
 M_A(S)=\mathcal S_{\theta,A}(M_B(S))
\end{equation}
is proved with all Gaussian prior cross terms in
Appendix~\ref{app:separator-hierarchy}. It is not an identification of the
smaller matrix with a principal block. Concavity gives
\[
 \sum_S\alpha_SM_A(S)\preceq
 \mathcal S_{\theta,A}\left(\sum_S\alpha_SM_B(S)\right).
\]
Apply monotonicity of the remaining parameter Schur map and of log
determinant. Successive Schur elimination equals joint elimination, so this
bounds the objective for $A$ by that for $B$ at the same $\alpha$.
Maximize over the common simplex. A point mass on any feasible schedule and
\eqref{eq:separator-schur} prove the first inequality.
\end{proof}

The empty-anchor level is the hull of the true information matrices. It can
still exceed the best single schedule: with $R=I_2$, $F=2I_2$,
$J_0=I_2$, and $k=1$, the two information matrices are
$\diag(5,1)$ and $\diag(1,5)$, whose equal mixture has determinant $9$
whereas each feasible determinant is $5$. This gap does not arise from
covariance truncation.

When every block is a singleton, mixture matrices depend only on the
marginal selection weights $z$. Their parameter Schur complement equals
\eqref{eq:vn-resolvent} for the diagonal residual split $D_A$; a direct
factor proof is in Appendix~\ref{app:separator-hierarchy}. With all times
anchored, $D_A=rI$, recovering the physical-nugget Liu split. Under the
right-endpoint convention and $b=1$, the final time is unanchored: for
stationary variance $P$ and $n\ge2$ the actual split is
$\diag(r,\ldots,r,r+P(1-\rho^2))$. For $n=1$ it is $[r+P]$.
This distinction matters when choosing a baseline. Improving on the physical
split does not prove dominance over every admissible virtual-noise split.

The statistical and optimization ingredients have direct predecessors.
Focused Bayesian design already uses marginalized nuisance precision
\citep{Alexanderian2021}; \citet[Section 7]{LevineHow2013} uses latent
augmentation to bound focused sensor information through an augmented-target
criterion. General PSD experiment atoms and their design hulls are established
\citep{HarmanTrnovska2009,Filova2026}. In particular, set
$E=[I_p,0]^T$ in the subsystem criterion of
\citet[Theorem 4.3 and Corollaries 4.4--4.5]{Sagnol2015}:
$(E^TM^{-1}E)^{-1}=\mathcal S_\theta(M)$.
Treating separator patterns and the fixed prior as atoms embeds this exact
integer design in their general conic framework. Our specialized result is
the comparison across the changing latent representations in
\eqref{eq:separator-cross-identity}--\eqref{eq:separator-order}, together
with short-block support certificates for the original correlated experiment.
Grouping discrete alternatives before convexification has a separate
disjunctive-programming precedent \citep{Balas1985,Papageorgiou2025}. That
principle explains a strength--enumeration tradeoff, but does not by itself
compare statistical representations whose latent variables change; the
cross-representation Schur identity supplies that step here.
To the best of our knowledge, this particular nested Markov-anchor hierarchy
is not stated in the inspected predecessors. This qualified claim does not
assert a new Schur complement, generalized D criterion, generic conic
formulation, or design-mixture algorithm. Its practical cost includes pattern
construction and the growing nuisance system, in addition to pricing.
